\chapter{KESIMPULAN}
\label{chap:kesimpulan}

\section{Jawaban Pertanyaan Kesimpulan}
Pertanyaan kesimpulan pada spesifikasi tugas didasarkan pada seluruh eksperimen. \textit{Task} 1 sampai 3, \textit{Task} 4 sisi FP-Growth, dan \textit{Task} 5 sudah dikerjakan sehingga pertanyaan 1 sampai 4 dapat dijawab. Jawaban pertanyaan 5 sampai 7 akan dilengkapi setelah \textit{Task} 6 dan Bagian E selesai.

\begin{enumerate}
\item \textbf{Konferensi dengan popularitas publikasi tertinggi.} NeurIPS, dengan 25.455 makalah (1,27\% dari 1.997.244 publikasi pada 2015--2025) dan 57.800 penulis unik. NeurIPS tetap berada di peringkat pertama ketika \textit{workshop} dan acara satelit dikeluarkan dari perhitungan. Peringkat berikutnya, yaitu CVPR dan ICASSP, bergantung pada definisi: CVPR turun ke peringkat enam jika \textit{CVPR Workshops} dikeluarkan.
\item \textbf{Asosiasi konferensi paling populer dengan konferensi lain.} NeurIPS memiliki asosiasi paling banyak, tetapi bukan yang paling kuat. NeurIPS termasuk dalam 42 pasangan \textit{frequent}, 35 dari 40 \textit{itemset} berisi lima atau enam konferensi, dan menjadi tujuan 240 \textit{rules}, terbanyak dari semua konferensi. Namun, median \textit{lift} \textit{rules} yang menuju NeurIPS hanya 23,61, kedua terendah. Semakin populer sebuah konferensi, semakin rendah \textit{lift} \textit{rules} yang menuju kepadanya (Spearman $-0{,}933$), karena \textit{lift} dibagi dengan \textit{support} konferensi tujuan. Asosiasi satu-ke-satu terkuat justru dimiliki pasangan yang lebih kecil, yaitu ECOC $\Rightarrow$ OFC (\textit{lift} 80,12) dan ASRU $\Rightarrow$ INTERSPEECH (74,17) (Subbab~\ref{sec:hasil-t5-lift}).
\item \textbf{Kombinasi konferensi yang paling sering muncul bersama.} ICML dan NeurIPS (18.503 penulis, \textit{support} 1,02\%), satu-satunya pasangan yang masih \textit{frequent} pada \textit{min\_sup} 1\%, disusul ICLR dan NeurIPS (17.852) serta AAAI dan NeurIPS (15.973). \textit{Triple} teratas adalah ICLR, ICML, dan NeurIPS (10.413 penulis), dan \textit{itemset} terpanjang berisi enam konferensi, yaitu AAAI, CVPR, ECCV, ICCV, ICLR, dan NeurIPS (2.005 penulis) (Subbab~\ref{sec:hasil-t4-fp-kombinasi}).
\item \textbf{\textit{Rule} dengan \textit{lift} tertinggi.} ECCV + ACM Multimedia + NeurIPS $\Rightarrow$ AAAI + CVPR + ICCV (\textit{lift} 172,27). Dari 2.983 penulis yang menerbitkan di ECCV, ACM Multimedia, dan NeurIPS, 61,2\% juga menerbitkan di AAAI, CVPR, dan ICCV, atau 172 kali lebih sering dibanding penulis pada umumnya. \textit{Rule} ini hanya didukung 1.827 penulis, sedikit di atas ambang 1.822, dan tidak lagi lolos jika profil nama umum yang menggabungkan banyak orang dikeluarkan sehingga perlu dibaca hati-hati. \textit{Lift} tingginya terutama berasal dari \textit{consequent} tiga konferensi yang jarang muncul bersama. Di antara \textit{rule} satu-ke-satu, \textit{lift} tertinggi dimiliki ECOC $\Rightarrow$ OFC (80,12): 66,5\% penulis ECOC juga menerbitkan di OFC (Subbab~\ref{sec:hasil-t5-top10} dan \ref{sec:hasil-t5-menarik}).
\item \textbf{Perbedaan pola pada \textit{author-based} dan \textit{author-year-based}.} Sebagian sudah dapat dijawab dari pratinjau pasangan. Kedua representasi menjawab pertanyaan yang berbeda. \textit{Support} pasangan pada \textit{author-year-based} median 0,29 kali \textit{author-based}, sekitar 23\% pasangan tidak pernah muncul, dan konferensi dua tahunan berselang seperti ICCV dan ECCV tidak dapat berpasangan sama sekali (10.051 transaksi pada \textit{author-based}, nol pada \textit{author-year-based}). Urutan pasangan teratas relatif terjaga (18 dari 20 sama, korelasi peringkat 0,892). Jawaban lengkap menunggu hasil \textit{mining}.
\item \textbf{Algoritma dengan performa terbaik.} Belum dapat dijawab (\textit{Task} 6).
\item \textbf{Pengetahuan atau \textit{insight} utama.} Sementara ini, popularitas terkonsentrasi dan makin terkonsentrasi pada konferensi kecerdasan buatan: porsi sepuluh konferensi teratas naik dari 7,0\% pada 2015 menjadi 11,7\% pada 2024, dan NeurIPS tumbuh 10,2 kali. Sebagian kecil penulis menanggung sebagian besar aktivitas (8,3\% penulis menyumbang 52,6\% pasangan makalah-penulis), dan pilihan representasi transaksi menentukan jenis pola yang dapat ditemukan. Kesimpulan akhir menunggu hasil \textit{mining}.
\end{enumerate}

\section{Keterbatasan dan Asumsi}
\begin{enumerate}
\item \textit{Venue} diambil dari segmen kedua \textit{key} sehingga \textit{workshop} dan acara pendamping ikut terhitung sebagai konferensi induknya (7,2\% makalah menurut \textit{booktitle}). Dampaknya besar pada CVPR, CHI, dan ICC (Tabel~\ref{tab:sensitivitas}).
\item \textit{Record} \texttt{inproceedings} yang \textit{key}-nya bukan \texttt{conf/} (3.161 makalah, 0,16\%) dibuang, termasuk beberapa prosiding \textit{workshop} kecil.
\item Data 2025 tidak lengkap karena konferensi dua tahunan genap dan konferensi tahunan yang belum terindeks sehingga tidak dipakai menghitung tren dan membuat transaksi \textit{author-year-based} pada periode terakhir bias.
\item Nama penulis umum tanpa sufiks disambiguasi dapat mewakili banyak orang (\textit{Yang Liu} tercatat pada 332 \textit{venue}), dan 7.573 kelompok nama yang hanya berbeda diakritik atau huruf besar-kecil tidak digabung. Akibatnya, penulis unik cenderung terhitung lebih rendah dan sebagian transaksi menjadi terlalu besar.
\item Aturan duplikat berdasarkan judul, tahun, dan \textit{venue} dapat membuang sebagian makalah berbeda yang berjudul generik (paling banyak 421 baris, 0,02\% data).
\item ICLR tidak tercatat pada DBLP untuk 2015 dan 2016 sehingga dikeluarkan dari \textit{heatmap} indeks.
\item Pengenalan acara satelit lewat kata kunci \textit{booktitle} adalah pendekatan. Beberapa acara resmi, misalnya \textit{CHI Extended Abstracts}, dapat dianggap bagian konferensi utama.
\item Model transaksi berupa himpunan tidak menyimpan jumlah makalah per \textit{venue}: 38,3\% makalah pada \textit{author-based} dan 17,4\% pada \textit{author-year-based} tidak menambah item baru (Subbab~\ref{sec:hasil-npaper}).
\item Pratinjau perbandingan pasangan hanya mencakup \textit{venue} dengan \textit{support} minimal 0,1\% pada \textit{author-based} (574 \textit{venue}) dan belum merupakan hasil \textit{mining}.
\item Popularitas hanya diukur dari aktivitas publikasi, bukan dari sitasi atau kualitas. Pengelompokan bidang konferensi pada narasi adalah pengetahuan umum, bukan atribut dari data DBLP. Pada Subbab~\ref{sec:hasil-t4} dan \ref{sec:hasil-t5}, nama bidang dicocokkan dengan judul prosiding pada data DBLP.
\item Pada \textit{min\_sup} 0,1\%, 291 transaksi dengan 100 konferensi atau lebih, yang hampir seluruhnya nama umum tanpa nomor disambiguasi, menentukan lolos tidaknya sebagian \textit{itemset}. Tanpa transaksi tersebut, 109 dari 1.514 \textit{frequent itemsets} (7\%) dan tujuh dari sepuluh \textit{rule} dengan \textit{lift} tertinggi tidak lagi lolos ambang. Hasil tetap dilaporkan apa adanya.
\item Dengan \textit{min\_conf} 50\% dan \textit{support} setiap konferensi di bawah 3,2\%, setiap \textit{rule} otomatis memiliki \textit{lift} minimal sekitar 15,8. \textit{Lift} pada data ini dipakai untuk membandingkan \textit{rules}, bukan sebagai ukuran mutlak.
\item \textit{Rules} hanya ada untuk kombinasi yang lolos \textit{min\_sup}. Hubungan kuat antara dua konferensi kecil tidak terlihat.
\end{enumerate}
